% !TEX root = ../main.tex
\section{Einleitung}

\subsection{Ausgangslage}

Selbstgeführte Wechselrichter am Netz kommen in rückspeisefähigen Frequenzumrichtern (netzseitiger Teil), in Blindleistungskompensatoren und in Photovoltaik-Wechselrichtern vor. Die Versuchsanleitung (A.~Omlin, HSLU T\&A, Stand 08.11.2024) vergleicht drei Varianten für den netzseitigen Teil eines Frequenzumrichters (Tabelle~\ref{tab:varianten}).

\begin{table}[htbp]
  \centering
  \caption{Varianten des netzseitigen Teils eines Frequenzumrichters nach Versuchsanleitung.}
  \label{tab:varianten}
  \small
  \begin{tabular}{@{}p{3.6cm}p{3.6cm}p{3.2cm}p{4.2cm}@{}}
    \toprule
    Variante & Energiefluss & $\cos\varphi$ & Oberschwingungen \\
    \midrule
    Diodenbrücke & nur Netz $\to$ Last, kein Rückspeisen & -- & -- \\
    Umkehrgleichrichter (Thyristor) & beide Richtungen & vom Steuerwinkel abhängig, im Normalbetrieb deutlich $<1$ & durch die Schaltung weitgehend gegeben \\
    Selbstgeführter Wechselrichter (dieser Versuch) & beide Richtungen & frei einstellbar, z.\,B.\ $\cos\varphi=1$ & durch höhere Taktfrequenz und/oder geschickte Pulsmuster deutlich reduzierbar \\
    \bottomrule
  \end{tabular}
\end{table}

\subsection{Ziel des Versuchs}

Im Versuch werden das Funktionsprinzip sowie typische Strom- und Spannungsverläufe eines einphasigen, selbstgeführten Wechselrichters am Netz praktisch untersucht. Es soll gezeigt werden, wie sich Wirk- und Blindleistung einstellen lassen und wie sich Pulsmuster auf die Oberschwingungen auswirken. Die Messungen fanden am 13.11.2026 und am 20.11.2026 statt.

Gemäss Versuchsanleitung sollen die Studierenden

\begin{itemize}
  \item das theoretisch behandelte Verhalten des selbstgeführten einphasigen Stromrichters überprüfen,
  \item wissen, wie sich Wirkleistung und Blindleistung einstellen lassen,
  \item wissen, wie sich verschiedene Pulsmuster auf die Oberschwingungen auswirken,
  \item Spannungs- und Stromverläufe bei selbstgeführten Stromrichtern sicher interpretieren können.
\end{itemize}

\subsection{Aufbau des Berichts}

Der Bericht folgt dem Aufgabenraster der Versuchsanleitung. Kapitel~\ref{sec:aufbau} beschreibt Versuchsaufbau, Messpunkte und die verwendete Vorzeichenkonvention, Kapitel~\ref{sec:theorie} fasst die benötigte Theorie zusammen. Kapitel~\ref{sec:vorbereitung} behandelt die Vorbereitungsaufgabe (Aufgabe~5), Kapitel~\ref{sec:messungen} die Messaufgaben 6.1 bis 6.4. Abweichungen und Unstimmigkeiten werden in Kapitel~\ref{sec:diskussion} diskutiert, das Fazit steht in Kapitel~\ref{sec:fazit}.
